PID controllers still regulate the large majority of loops in the chemical process industries \cite{astrom2006pid, seborg2010process}, and the temperature loop of an exothermic continuous stirred tank reactor (CSTR) is one of the most demanding among them. Arrhenius kinetics couple reactor temperature and reaction rate through an exponential positive feedback, which produces steady-state multiplicity, parametric sensitivity and open-loop unstable regimes \cite{uppal1974dynamic, aris1958analysis, bequette2003process}. The cooling valve that closes the loop is bounded in stroke and in travel speed, and when an aggressive controller drives it into saturation the loop is effectively opened and the integrator winds up \cite{astrom1989windup, hippe2006windup, galeani2009tutorial, tarbouriech2009antiwindup}.

The tuning rules most often used as references for this loop are the reaction-curve rules of Ziegler and Nichols (1942) \cite{ziegler1942optimum} and Cohen and Coon (1953) \cite{cohen1953theoretical}, derived for linear, unconstrained processes \cite{odwyer2009handbook, shinskey1996process}. Their failure on exothermic reactors is a recurring finding \cite{visioli2006practical, luyben2007chemical}, and it is frequently used to motivate model-based, nonlinear or learning-based alternatives. This paper is not another ranking of tuning rules. It is an implementation study that asks how much of such verdicts is decided by details of the implementation rather than by the rule: the actuator model, the anti-windup logic and the measurement path. When the anti-windup logic is wrong the resulting windup is attributed to the tuning rule; when the measurement is unfiltered, the noise vulnerability is attributed to a sensitivity peak that the physical plant does not have.

\subsection{Positioning against the anti-windup literature}
\label{subsec:positioning}
Anti-windup for PID loops is a mature topic with four families of solutions (Table~\ref{tab:aw_families}). Conditional integration (clamping) and back-calculation are the schemes implemented in practice \cite{astrom1989windup, hippe2006windup, peng1996nonlinear, bohn1995analysis}, and they are the two compared here. Compensators designed by linear matrix inequalities treat magnitude and rate saturation together \cite{galeani2009tutorial, tarbouriech2009antiwindup, lerch2021antiwindup}; they are not part of the comparison because they require a plant model and a design step that the PID practitioner studied here does not have, and the question asked is what an implementation of the two industrial schemes does. Recent work refines individual schemes: rate-limited anti-windup for a fractional-order PID \cite{wang2020rate}, gradient-based optimization of anti-windup PID gains \cite{aksland2022gradient}, and learned or operator-theoretic variants of back-calculation \cite{serrano2024neuron, morohoshi2025nonlinear}. Two facts frame the present study. First, conditional integration is correctly specified by the direction of the integrator relative to the saturation (freeze when the increment pushes the command further into the stop), but it is easily implemented as the shortcut ``freeze when the command exceeds the stop and the error is positive'', which is right for a direct-acting controller only. In a reverse-acting loop ($K_p < 0$) the sign of the error that pushes the command into the stop is reversed, and a shortcut that keeps the direct-acting convention is \emph{sign-blind}. We call it \emph{sign-blind amplitude clamping} and present it as a frequent implementation error, not as a rule attributed to any publication. Second, back-calculation depends on a tracking time constant $T_t$ that is rarely reported together with the tuning rule it accompanies.

\begin{table}[t]
\centering
\caption{Anti-windup families and what each assumes about the loop. The last row is the scheme adopted in this study; it is the sign-of-increment form of conditional integration \cite{peng1996nonlinear, hippe2006windup} evaluated on the slew-limited valve position, not a new algorithm.}\label{tab:aw_families}
\footnotesize
\begin{tabular}{p{2.6cm}p{2.6cm}p{2.0cm}p{1.6cm}p{1.4cm}}
\toprule
Scheme & Integrator logic & Loop sign & Rate limit & Extra parameter \\
\midrule
Sign-blind amplitude clamping (implementation error) & freeze if $u_{\text{raw}}$ is beyond a stop and $e$ has the saturating sign of a direct-acting loop & direct-acting convention reused & not seen & none \\
Back-calculation \cite{astrom1989windup, hippe2006windup} & integrate $(u_{\text{act}} - u_{\text{raw}})/T_t$ & independent & seen through $u_{\text{act}}$ & $T_t$ \\
Dynamic compensators \cite{galeani2009tutorial, tarbouriech2009antiwindup, lerch2021antiwindup} & LMI-designed filter on the saturation error & through the model & designed for & model, LMI weights \\
Directional rate-aware clamping (Eq.~\eqref{eq:clamping_antiwindup}) & freeze iff $(u_{\text{raw}}-u_{\text{act}})\,\delta I > 0$ & independent (product form) & seen through $u_{\text{act}}$ & none (but depends on $S_{\max}$, $\Delta t$) \\
\bottomrule
\end{tabular}
\end{table}

What the literature does not provide, to our knowledge, is a controlled comparison of these schemes on a reverse-acting thermal loop with a slew-limited valve in which the tuning rule is held fixed, together with the mechanism by which the sign-blind rule fails. That is the gap addressed here. The contribution is an analysis of implementation, not a new anti-windup law.

\subsection{Study design and contributions}
The study uses a three-state CSTR model with an explicit jacket balance and a cooling valve limited to $[0, 300]$~L/min and $80$~L/min$^2$. Eight PID tunings, including modern robust baselines \cite{tyreus1992tuning, astrom2004revisiting, skogestad2003simple, grimholt2018optimal, garpinger2014performance, padula2011tuning}, are run under one protocol with error metrics evaluated on the true reactor temperature. The contributions are:
\begin{enumerate}
    \item A seven-policy ablation of the integrator logic (no anti-windup, back-calculation with $T_t$ scaled by $0.1$ and $10$, sign-blind amplitude clamping, and the $2\times 2$ factorial of trigger---amplitude or amplitude plus slew rate---and sign awareness---unconditional or directional), run on three setpoint steps and a load profile, with a mechanistic account of why sign-blind clamping latches a reverse-acting loop at full cooling. A proposition shows why the sign-blind rule latches; the slew-rate trigger, not sign awareness, is what restores the loop, and the resulting rule needs no per-loop constant but depends on the slew limit, which is mapped over $20$--$320$~L/min$^2$; back-calculation is scanned over its tracking constant and the ablation is repeated at a second operating point.
    \item A quantification of what the design-model choice costs: the $M_s \le 1.6$ bound is imposed on the identified dead-time model (the proxy used by most tuning studies \cite{garpinger2015software, sato2019optimal}) and on the exact linearization, and the setpoint weights are optimized jointly with the gains.
    \item A benchmark of eight tunings covering servo tracking, out-of-sample steps, unmeasured load disturbances (deterministic and replicated) and robustness on the exact linearized plant, including $\|CS\|_\infty$ and the high-frequency controller gain.
    \item A noise campaign with and without a first-order measurement filter, and a study of the derivative-filter coefficient $N$, so that the noise-to-valve mechanism is attributed to the controller gain at high frequency rather than to a sensitivity peak \cite{segovia2013noise}.
    \item A continuous map of the trade-off between tracking error and valve travel \cite{alcantara2013pid, marques2020applications}, and a global sensitivity analysis with Sobol' indices \cite{sobol2001global, saltelli2010variance} with bootstrap intervals, a stated starting point and a direct test of the heat-removal-margin explanation.
    \item A pseudocode of the integrator update and a five-item checklist for reporting the anti-windup logic of a PID comparison (Section~\ref{subsec:checklist}).
\end{enumerate}
Section~\ref{sec:model} presents the plant and the actuator limits, Section~\ref{sec:pid} the controller structure, the anti-windup variants and the tuning rules, Section~\ref{sec:results} the results and Section~\ref{sec:conclusion} the conclusions.
